%%%%%%%%%%%%%%%%%%%%%%%%%%%%%%%%%%%%%%%%%%%%%%%%%%%%%%%%%%%%%%%%%%%%%%%%%%
% Symbol definitions
%%%%%%%%%%%%%%%%%%%%%%%%%%%%%%%%%%%%%%%%%%%%%%%%%%%%%%%%%%%%%%%%%%%%%%%%%%

\glsxtrnewsymbol[description={Category of contexts (\textsf{FinSet})}]{CatOfCtx}{\ensuremath{\mathbb{F}}}

\glsxtrnewsymbol[description={Continuation Monad $R^{R^{-}}$ on ${\textsf{Set}}^{\glspt{CatOfCtx}}$ with $R \in {\textsf{Set}}^{\glspt{CatOfCtx}}$}]{ContMonad}{\ensuremath{K_R}}

\glsxtrnewsymbol[description={Variable presheaf $Y(1) \in {\textsf{Set}}^{\glspt{CatOfCtx}}$}]{VarPresheaf}{\ensuremath{L}}

\glsxtrnewsymbol[description={Label type}]{LabelType}{\ensuremath{\ell}}

\glsxtrnewsymbol[description={Context},user1={\ensuremath{\Gamma}},user2={\ensuremath{\Delta}},user3={\ensuremath{\Psi}}]{Ctx}{\ensuremath{\Gamma, \Delta, \Psi}}

\glsxtrnewsymbol[description={Pure term},user1={\ensuremath{V}},user2={\ensuremath{W}}]{Pure}{\ensuremath{V, W}}

\glsxtrnewsymbol[description={Effectful term},user1={\ensuremath{M}},user2={\ensuremath{N}}]{Comp}{\ensuremath{M, N}}

\glsxtrnewsymbol[description={Variable},user1={\ensuremath{x}},user2={\ensuremath{y}},user3={\ensuremath{z}}]{Var}{\ensuremath{x, y, z}}

\glsxtrnewsymbol[description={Unit term}]{UnitTerm}{\ensuremath{\ast}}

\glsxtrnewsymbol[description={Lambda abstraction operator}]{LamOp}{\textsf{lam}}

\glsxtrnewsymbol[description={Pair operator}]{PairOp}{\textsf{pair}}

\glsxtrnewsymbol[description={Push operator}]{PushOp}{\textsf{push}}

\glsxtrnewsymbol[description={Return operator}]{ReturnOp}{\textsf{return}}

\glsxtrnewsymbol[description={Application operator}]{AppOp}{\textsf{app}}

\glsxtrnewsymbol[description={Jump operator}]{VarOp}{\textsf{var}}

\glsxtrnewsymbol[description={Branch operator}]{SubOp}{\textsf{sub}}

\glsxtrnewsymbol[description={Pattern match operator}]{PmOp}{\textsf{pm}}

\glsxtrnewsymbol[description={Type},user1={\ensuremath{X}},user2={\ensuremath{Y}},user3={\ensuremath{Z}}]{Type}{\ensuremath{X, Y, Z}}

\glsxtrnewsymbol[description={Membership judgement},user1={\ensuremath{\in}}]{Member}{\ensuremath{x : X \in \Gamma}}

\glsxtrnewsymbol[description={Typing judgement for pure terms},user1={\ensuremath{\vdash^v}}]{PureJudge}{\ensuremath{\Gamma \vdash^v V : X}}

\glsxtrnewsymbol[description={Typing judgement for computations},user1={\ensuremath{\vdash^c}}]{CompJudge}{\ensuremath{\Gamma \vdash^c M : X}}

\glsxtrnewsymbol[description={Appending a type to a context},user1={\ensuremath{\cdot}}]{CtxCons}{\ensuremath{\Gamma \cdot X}}

\glsxtrnewsymbol[description={Function type},user1={\ensuremath{\Rightarrow}}]{ArrowType}{\ensuremath{X \Rightarrow Y}}

\glsxtrnewsymbol[description={Unit type}]{UnitType}{\ensuremath{\mathbf{1}}}

\glsxtrnewsymbol[description={Product type},user1={\ensuremath{\times}}]{ProdType}{\ensuremath{X \times Y}}

%%%%%%%%%%%%%%%%%%%%%%%%%%%%%%%%%%%%%%%%%%%%%%%%%%%%%%%%%%%%%%%%%%%%%%%%%%
% Macros
%%%%%%%%%%%%%%%%%%%%%%%%%%%%%%%%%%%%%%%%%%%%%%%%%%%%%%%%%%%%%%%%%%%%%%%%%%

\newcommand\glsi[1]{\glshyperlink[\glsentryuseri{#1}]{#1}\glsadd{#1}\xspace}
\newcommand\glsii[1]{\glshyperlink[\glsentryuserii{#1}]{#1}\glsadd{#1}\xspace}
\newcommand\glsiii[1]{\glshyperlink[\glsentryuseriii{#1}]{#1}\glsadd{#1}\xspace}

%%%%%%%%%%%%%%%%%%%%%%%%%%%%%%%%%%%%%%%%%%%%%%%%%%%%%%%%%%%%%%%%%%%%%%%%%%
% Custom glossary style
%%%%%%%%%%%%%%%%%%%%%%%%%%%%%%%%%%%%%%%%%%%%%%%%%%%%%%%%%%%%%%%%%%%%%%%%%%

\newlength{\myglosssymbolwidth}\setlength{\myglosssymbolwidth}{0.15\textwidth}
\newlength{\myglossdescwidth}\setlength{\myglossdescwidth}{0.55\textwidth}
\newlength{\myglosskeywidth}\setlength{\myglosskeywidth}{0.2\textwidth}

\newglossarystyle{long-symbols}{%
% put the glossary in a longtable environment:
\renewenvironment{theglossary}%
%% To determine width of colums automatically: {\begin{longtable}{lp{\glsdescwidth}p{\glspagelistwidth}}}%
{\begin{longtable}{p{\myglosssymbolwidth}p{\myglossdescwidth}p{\myglosskeywidth}}}%
{\end{longtable}}%
% Set the table's header:
\renewcommand*{\glossaryheader}{%
\bfseries Symbol & \bfseries Description & \bfseries Key
\\\endhead}%
% No heading between groups:
\renewcommand*{\glsgroupheading}[1]{}%
\renewcommand*{\glossentry}[2]{%
\glstarget{##1}{\glossentryname{##1}}
& \glossentrydesc{##1}% Description
& \texttt{##1}%
\tabularnewline % end of row
}%
% Similarly for sub-entries (no sub-entry numbers)
\renewcommand*{\subglossentry}[3]{%
% ignoring first argument (sub-level)
\glstarget{##2}{\glossentryname{##2}}
& \glossentrydesc{##2}% Description
& \texttt{##2}%
\tabularnewline % end of row
}%
% Nothing between groups:
\renewcommand*{\glsgroupskip}{}%
}
